\section{The percolation sum in Hastings' cluster expansion}
\label{sec:hastings_percolation}

% Hastings, arXiv:cond-mat/0508554v3, proj3new.tex, lines 1196--1254 (the
% inclusion--exclusion sum over large clusters and eq. perc2).

At high temperature, \cite[lines 1196--1225]{Hastings2006Solving} bounds
$e^{-\beta H}-\rho(\beta,l_{\rm proj})$ via an inclusion--exclusion sum over
families $U$ of clusters of diameter $>l_{\rm proj}$, where each $U$ contributes
at most $\|O\|\prod_{C\in U}a_C$ with $0\le a_C\le y^{n_C}$ and
$y=(e^{J\beta}-1)e^{SJ\beta}$. For terms of range $R$, a cluster of $n$ terms
has diameter at most $2Rn$, so diameter $>l_{\rm proj}$ implies
$n_C\ge l+1$ with $l=\lfloor l_{\rm proj}/(2R)\rfloor$, correcting the
nearest-neighbour threshold $l_{\rm proj}+1$ of
\cite[line 1240]{Hastings2006Solving}; see
\cite{gap:hastings05_percolation_cluster_size}.

\begin{theorem}[Families of a given number of clusters]
    \label{thm:hastings_percolation_families}
    \lean{TNLean.Lattice.sum_powersetCard_prod_le}
    \leanok
    For a finite set $T$ with weights $a_C\ge0$ ($C\in T$) and $m\in\mathbb N$,
    \begin{align}
        \sum_{U\subseteq T,\ |U|=m}\ \prod_{C\in U}a_C
        \le\frac1{m!}\Bigl(\sum_{C\in T}a_C\Bigr)^m
        \label{eq:hastings_percolation_families}
    \end{align}
    \cite[line 1210]{Hastings2006Solving}.
\end{theorem}
\begin{proof}
    \leanok
    Induct on $|T|$: for $b\notin T$ and $s=\sum_{C\in T}a_C$,
    $s^{m+1}/(m+1)!+a_bs^m/m!\le(s+a_b)^{m+1}/(m+1)!$.
\end{proof}

\begin{theorem}[All nonempty families of clusters]
    \label{thm:hastings_percolation_exp}
    \lean{TNLean.Lattice.sum_nonempty_powerset_prod_eq,
        TNLean.Lattice.sum_nonempty_powerset_prod_le_exp}
    \leanok
    For a finite set $T$ with weights $a_C\ge0$ ($C\in T$),
    \begin{align}
        \sum_{\emptyset\ne U\subseteq T}\ \prod_{C\in U}a_C
        =\prod_{C\in T}(1+a_C)-1
        \le\exp\Bigl(\sum_{C\in T}a_C\Bigr)-1
        \label{eq:hastings_percolation_exp}
    \end{align}
    \cite[lines 1238--1241]{Hastings2006Solving}.
\end{theorem}
\begin{proof}
    \leanok
    Expand $\prod_{C\in T}(1+a_C)$ and use $1+a_C\le e^{a_C}$.
\end{proof}

\begin{theorem}[The sum over large clusters]
    \label{thm:hastings_percolation_cluster_sum}
    \lean{TNLean.Lattice.sum_pow_le_of_forall_lt,
        TNLean.Lattice.sum_pow_le_of_card_le}
    \leanok
    If every $C\in T$ has $n_C\ge l+1$, $|\{C\in T:n_C=k\}|\le Nx^k$ for all
    $k$ ($N\ge0$), and $x,y\ge0$ with $xy<1$, then
    \begin{align}
        \sum_{C\in T}y^{n_C}
        \le N\frac{(xy)^{l+1}}{1-xy}
        \label{eq:hastings_percolation_cluster_sum}
    \end{align}
    \cite[lines 1238--1241]{Hastings2006Solving}.
\end{theorem}
\begin{proof}
    \leanok
    Group by $k=n_C\ge l+1$ and bound
    \begin{align}
        \sum_{k\ge l+1}N(xy)^k=N\frac{(xy)^{l+1}}{1-xy}.
        \notag
    \end{align}
\end{proof}

\begin{theorem}[The percolation bound]
    \label{thm:hastings_percolation}
    \lean{TNLean.Lattice.sum_nonempty_powerset_prod_le_percolation,
        TNLean.Lattice.pow_succ_div_one_sub_le,
        TNLean.Lattice.pow_eq_exp_neg_div}
    \leanok
    \uses{thm:hastings_percolation_exp, thm:hastings_percolation_cluster_sum}
    In the setting of Theorem~\ref{thm:hastings_percolation_cluster_sum}, if
    $0\le a_C\le y^{n_C}$, $0<xy\le1/2$, and $\xi=1/\ln(1/(xy))$, then
    \begin{align}
        \sum_{\emptyset\ne U\subseteq T}\ \prod_{C\in U}a_C
        \le\exp\bigl(Ne^{-l/\xi}\bigr)-1
        \label{eq:hastings_percolation}
    \end{align}
    \cite[eq.~(perc2), lines 1227--1254]{Hastings2006Solving},
    \cite{gap:hastings05_percolation_cluster_size}.
\end{theorem}
\begin{proof}
    \leanok
    \uses{thm:hastings_percolation_exp, thm:hastings_percolation_cluster_sum}
    Combine (\ref{eq:hastings_percolation_exp}) and
    (\ref{eq:hastings_percolation_cluster_sum}) with
    $(xy)^{l+1}/(1-xy)\le(xy)^l=e^{-l/\xi}$ for $0<xy\le1/2$.
\end{proof}

\begin{theorem}[The percolation bound for connected clusters]
    \label{thm:hastings_percolation_connected}
    \lean{TNLean.Lattice.sum_nonempty_powerset_prod_le_percolation_connected}
    \leanok
    \uses{thm:hastings_percolation, thm:hastings_connected_set_count}
    Let $G$ be a simple graph on $N$ vertices of maximum degree $\Delta$,
    $x=\max(1,\Delta^2)$, and $T$ be a family of connected vertex sets $C$ with
    $|C|\ge l+1$ and weights $0\le a_C\le y^{|C|}$. If $0<xy\le1/2$, then
    (\ref{eq:hastings_percolation}) holds with $\xi=1/\ln(1/(xy))$.
\end{theorem}
\begin{proof}
    \leanok
    \uses{thm:hastings_percolation, thm:hastings_connected_set_count}
    Combine Theorem~\ref{thm:hastings_connected_set_count}
    ($|\{C\in T:|C|=k\}|\le Nx^k$) with Theorem~\ref{thm:hastings_percolation}.
\end{proof}

\begin{theorem}[A percolation temperature]
    \label{thm:hastings_percolation_temperature}
    \lean{TNLean.Lattice.exists_percolation_temperature}
    \leanok
    For any $J,S,x\in\mathbb R$ there exists $\beta_0>0$ such that
    $x(e^{J\beta}-1)e^{SJ\beta}\le1/2$ for all $0\le\beta\le\beta_0$
    \cite[lines 1242--1245]{Hastings2006Solving}.
\end{theorem}
\begin{proof}
    \leanok
    The function $\beta\mapsto x(e^{J\beta}-1)e^{SJ\beta}$ is continuous and
    vanishes at $\beta=0$.
\end{proof}

\begin{definition}[Clusters of large diameter]
    \label{def:hastings_large_diameter}
    \lean{TNLean.Lattice.ClusterExpansion.IsLargeDiam}
    \leanok
    \uses{def:hastings_cluster_of_terms}
    For terms $H_k$ supported in a metric space of sites, a set $C$ of terms has
    \emph{large diameter} if the support of $H_C$ contains two sites at distance
    $>l_{\mathrm{proj}}$ \cite[lines 1104--1107]{Hastings2006Solving}.
\end{definition}

\begin{theorem}[Size and number of clusters of large diameter]
    \label{thm:hastings_large_diameter_count}
    \lean{TNLean.Lattice.ClusterExpansion.dist_le_of_walk,
        TNLean.Lattice.ClusterExpansion.dist_le_of_mem_biUnion_of_connected,
        TNLean.Lattice.ClusterExpansion.card_ge_of_isLargeDiam,
        TNLean.Lattice.ClusterExpansion.card_add_card_boundary_le}
    \leanok
    \uses{def:hastings_large_diameter}
    If every term support has diameter at most $2R$ ($R>0$) and overlaps at most
    $\Delta$ other terms, then any connected set of $n$ terms has support
    diameter at most $2Rn$, every large-diameter connected set $C$ satisfies
    $n_C\ge\lfloor l_{\mathrm{proj}}/(2R)\rfloor+1$, and
    $n_C+n_{\partial C}\le(\Delta+1)n_C$
    \cite[line 1235]{Hastings2006Solving},
    \cite{gap:hastings05_percolation_cluster_size}.
\end{theorem}
\begin{proof}
    \leanok
    Any two sites in the support of $C$ are connected by an overlap walk of at
    most $n$ supports of diameter $\le2R$, giving distance $\le2Rn$, and each
    $k\in\partial C$ overlaps some $c\in C$.
\end{proof}

\begin{theorem}[The percolation bound for clusters of large diameter]
    \label{thm:hastings_percolation_large_diameter}
    \lean{TNLean.Lattice.ClusterExpansion.norm_trace_mul_exp_sub_truncatedExp_le_isLargeDiam,
        TNLean.Lattice.ClusterExpansion.exp_neg_floor_div_le,
        TNLean.Lattice.ClusterExpansion.norm_trace_mul_exp_sub_truncatedExp_le_isLargeDiam_dist}
    \leanok
    \uses{thm:hastings_large_diameter_count, thm:hastings_truncation_trace_bound,
        thm:hastings_percolation_connected}
    In the setting of Theorem~\ref{thm:hastings_large_diameter_count}, let
    $H=\sum_kH_k$ have $N$ Hermitian terms with $\|H_k\|\le J$, $\beta\ge0$,
    $x=\max(1,\Delta^2)$, $y=(e^{J\beta}-1)e^{(\Delta+1)J\beta}$, and
    $0<xy\le1/2$. With $Z=\operatorname{tr}e^{-\beta H}$,
    $l=\lfloor l_{\mathrm{proj}}/(2R)\rfloor$, $\xi=1/\ln(1/(xy))$, and
    $\xi_{\mathrm{perc}}=2R/\ln(1/(xy))$, for every $O$,
    \begin{align}
        \bigl|\operatorname{tr}\bigl(O(e^{-\beta H}-\rho(\beta,l_{\mathrm{proj}}))\bigr)\bigr|
         & \le\|O\|\bigl(\exp(Ne^{-l/\xi})-1\bigr)Z
        \notag                                                                               \\
         & \le\|O\|\bigl(\exp(N(xy)^{-1}e^{-l_{\mathrm{proj}}/\xi_{\mathrm{perc}}})-1\bigr)Z
        \label{eq:hastings_percolation_large_diameter}
    \end{align}
    \cite[eq.~(perc2), lines 1227--1254]{Hastings2006Solving},
    \cite{gap:hastings05_percolation_cluster_size,gap:hastings05_cluster_boundary}.
\end{theorem}
\begin{proof}
    \leanok
    \uses{thm:hastings_large_diameter_count, thm:hastings_truncation_trace_bound,
        thm:hastings_percolation_connected}
    Combine Theorem~\ref{thm:hastings_truncation_trace_bound} ($S=\Delta+1$),
    Theorem~\ref{thm:hastings_large_diameter_count}, and
    Theorem~\ref{thm:hastings_percolation_connected}, using
    $l\ge l_{\mathrm{proj}}/(2R)-1$.
\end{proof}
